\documentclass[11pt]{article}

\usepackage[a4paper,margin=2.35cm]{geometry}
\usepackage{fontspec}
\setmainfont{DejaVu Serif}
\setsansfont{DejaVu Sans}
\setmonofont{DejaVu Sans Mono}[Scale=0.84]
\usepackage{microtype}
\usepackage{booktabs}
\usepackage{longtable}
\usepackage{tabularx}
\usepackage{array}
\usepackage{amsmath,amssymb}
\usepackage{xcolor}
\usepackage{enumitem}
\usepackage{graphicx}
\usepackage{caption}
\usepackage{float}
\usepackage{fancyhdr}
\usepackage{seqsplit}
\usepackage{url}
\usepackage[hidelinks,unicode]{hyperref}

% Shared immutable facts for the English and Chinese manuscripts.
% Values are transcribed from content-addressed JSON artifacts.
\newcommand{\PaperDate}{16 July 2026}
\newcommand{\DevCases}{9}
\newcommand{\DevSemanticCases}{2}
\newcommand{\QualificationCases}{8}
\newcommand{\QualificationTargets}{74}
\newcommand{\FocusedTests}{53}
\newcommand{\FullTests}{332}

\newcommand{\AccA}{0.976608}
\newcommand{\CovA}{0.988304}
\newcommand{\RiskA}{0.011834}
\newcommand{\AurcA}{0.014603}
\newcommand{\AccB}{0.976608}
\newcommand{\CovB}{0.988304}
\newcommand{\RiskB}{0.011834}
\newcommand{\AurcB}{0.014603}
\newcommand{\AccC}{0.976608}
\newcommand{\CovC}{0.988304}
\newcommand{\RiskC}{0.011834}
\newcommand{\AurcC}{0.014603}
\newcommand{\AccD}{0.988304}
\newcommand{\CovD}{1.000000}
\newcommand{\RiskD}{0.011696}
\newcommand{\AurcD}{0.014705}

\newcommand{\DevBCalls}{2}
\newcommand{\DevBInput}{5752}
\newcommand{\DevBOutput}{381}
\newcommand{\DevBLatency}{4128.415}
\newcommand{\DevDCalls}{6}
\newcommand{\DevDInput}{7538}
\newcommand{\DevDOutput}{3598}
\newcommand{\DevDLatency}{22171.453}

\newcommand{\QualBCalls}{16}
\newcommand{\QualBInput}{76976}
\newcommand{\QualBOutput}{26436}
\newcommand{\QualBMeanLatency}{27.95}
\newcommand{\QualDCalls}{59}
\newcommand{\QualDInput}{79571}
\newcommand{\QualDOutput}{19547}
\newcommand{\QualDMeanLatency}{44.61}
\newcommand{\QualMaxObservedContext}{13306}
\newcommand{\QualMaxBudgetContext}{14142}
\newcommand{\LoadedContext}{16384}

\newcommand{\DevReportHash}{1e77f66d924622153798743c209497b5025271c5d94ad4cebf6290871ff87e27}
\newcommand{\DevManifestHash}{a1a199272b9019dd99dab873239f18ba1d8b656235c54ac40ad84fb5d1738797}
\newcommand{\QualManifestHash}{71aa69e25fddf8a3cfe9ad78e2cc2ba08015347f8ea53166985c081120fe2666}
\newcommand{\VocabularyHash}{874548a10f8871449be9b749c99e95e72e8852d261324abc8fa3c492246ed888}
\newcommand{\InvalidQualHash}{d0a5fb9a82d41bda557669fd384820485c5265f187679f2e954ad347eec25cde}
\newcommand{\ValidQualHash}{c89145cf23943fd2991504a31cea5a10ceddccc90725e31932cc02770b1c0538}
\newcommand{\QualAssessmentHash}{bec8279af416b796d4ea5fe071a8947fcd1d7d123c1cc021168dd928fcb87efb}
\newcommand{\DatasetPromptHash}{1b75583a35549a5cbdfecc79be5796eb358d2e5f976f5483b25756a81f8f8060}
\newcommand{\DatasetSchemaHash}{e79212874554daf93c401578b97963f45934854b22b6786ee41ae60bf8dfd7ca}
\newcommand{\LegacyPromptHash}{3d4047f7bdc36e6ae832fe7e988bd06bff39be81f73bea101c6abdfe32b78a0f}
\newcommand{\LegacySchemaHash}{83ee4fb51579af04a304c61eb4b610263b1d02fcae1a90a56cfdb131e445a91f}


\definecolor{DeepBlue}{HTML}{174A7E}
\definecolor{SoftBlue}{HTML}{EAF2F8}
\definecolor{FreezeRed}{HTML}{9E2A2B}
\definecolor{FreezeBg}{HTML}{FFF3F3}
\definecolor{GrayText}{HTML}{555555}

\newcommand{\artifacthash}[1]{{\scriptsize\ttfamily\expandafter\seqsplit\expandafter{#1}}}
\newcommand{\status}[1]{\textsf{\MakeUppercase{#1}}}
\newcommand{\NA}{\textsf{N/A}}
\newcommand{\placeholder}{\texttt{<TO BE BOUND BEFORE GOLD UNSEALING>}}
\newcolumntype{Y}{>{\raggedright\arraybackslash}X}
\newcolumntype{C}{>{\centering\arraybackslash}X}
\newcolumntype{R}{>{\raggedleft\arraybackslash}p{1.75cm}}
\setlength{\emergencystretch}{1.5em}

\pagestyle{fancy}
\fancyhf{}
\fancyhead[L]{Evidence-Constrained Semantic Schema Architecture}
\fancyhead[R]{Pre-blind technical report}
\fancyfoot[C]{\thepage}
\setlength{\headheight}{14pt}

\title{\textbf{How Much Semantic Architecture Is Necessary?}\\
\large A Causally Controlled A/B/C/D Protocol for Evidence-Constrained Dataset Schema Induction}
\author{Anonymous pre-blind technical report\\
\small Authors and affiliations to be bound at protocol freeze}
\date{\PaperDate}

\begin{document}
\maketitle

\begin{abstract}
This work studies the semantic component of a system whose product interface is
\emph{dataset in, deterministic schema out}. The physical schema remains parser-owned;
language models may only propose property-level semantic additions over existing
fields and cited evidence. We replace a historically fragmented semantic pipeline
with a controlled four-condition experiment: A is deterministic only; B adds at
most one dataset-level semantic call; C reuses the exact B response and applies
deterministic evidence-identity, field-relevance, and conflict checks; D preserves
the legacy per-field/manual-group schedule and non-empty-evidence merge behavior.

The contribution at the present stage is an auditable protocol and validated
measurement harness, not a claim that any architecture wins. On nine synthetic
development cases, only two exposed dataset-level semantic opportunities and one
opportunity was absent from gold. All pairwise architecture comparisons were
therefore correctly declared non-comparable. A smaller Qwen3.5 9B backend added no
accepted, scorable claim in this set, making the observed B--C difference
uninformative rather than a null effect. Adversarial tests established the intended
call, replay, evidence, abstention, applicability, unsupported-claim, and AURC
invariants.

Backend qualification subsequently exposed a runtime and measurement failure: an
8,192-token Qwen3.6-27B run with thinking effectively enabled was invalidated. After
repairing replay and telemetry provenance, disabling thinking, and loading a
16,384-token context, the same Q4\_K\_M model passed every predeclared operational
criterion over eight non-blind calibration datasets and 74 semantic targets. This
qualification establishes execution compatibility only. The causal contribution
of dataset-level reasoning and deterministic verification remains reserved for an
independently annotated, frozen blind evaluation. Explicit provenance placeholders
are included for that freeze.
\end{abstract}

\noindent\textbf{Keywords:} semantic schema induction; selective prediction;
evidence grounding; deterministic verification; replay identity; model capability
debt; dataset metadata.

\section{Introduction}

Dataset schema induction contains two different epistemic problems. Physical facts
such as field names, stored types, shapes, groups, and declared attributes are
properties of bytes and format standards. Contextual meaning---whether a column is
an identifier, a measurement, a time axis, or a domain quantity---often requires
document-level semantic integration. Treating both problems as free-form language
generation makes plausible metadata difficult to distinguish from schema truth.
Treating both as deterministic parsing leaves genuinely contextual information
unused.

The system studied here therefore retains a strict authority boundary. Deterministic
extractors own canonical physical structure. A semantic component may propose
claims only for existing fields, cite evidence identifiers from the input record,
and abstain when support is insufficient. This design is aligned with provenance
models that make derivation auditable \cite{prov2013}, evidence-backed model revision
\cite{gao2023rarr}, and selective prediction, where lower coverage can be exchanged
for lower risk \cite{geifman2017selective}.

The historical implementation used per-field calls, manual grouping, retries, and
multiple merge rules. Some of this complexity may have been useful when smaller
models struggled to integrate repository- or dataset-level context. We call
complexity that remains after its compensating model limitation has weakened
\emph{Model Capability Debt}. The scientific question is not whether a newer model
can be inserted everywhere. It is whether the minimum architecture that preserves
the trust boundary can be identified empirically.

We ask three causal questions:
\begin{enumerate}[label=\textbf{RQ\arabic*:},leftmargin=2.2cm]
  \item What does one dataset-level semantic reasoning call add beyond deterministic extraction?
  \item Holding the semantic response exactly constant, what does deterministic verification change?
  \item Does the minimal verified hybrid preserve trustworthiness relative to the legacy per-field system while using fewer semantic calls?
\end{enumerate}

At this pre-blind stage, none of these questions has a valid effect estimate. The
completed work establishes that the future estimates can be causally interpreted,
and identifies precisely what evidence is still missing.

\section{Problem Formulation and Product Boundary}

The product-level mapping is
\[
  \texttt{dataset file(s)} \longrightarrow
  \texttt{deterministic canonical schema} \oplus
  \texttt{optional semantic claim ledger}.
\]
The deterministic schema is a complete standalone output. Semantic augmentation is
an optional overlay and may never invent fields, silently overwrite stronger
deterministic claims, or convert missing support into confidence.

For field path $f$, property $p$, and proposed value $v$, a semantic claim is
represented as
\[
  c=(f,p,v,E,s,d),
\]
where $E$ is a set of evidence identifiers, $s$ is a confidence value used for
risk--coverage ordering, and $d$ is the final decision. Correctness and confidence
are deliberately separate: a claim may be confident but unsupported, or supported
but uncertain. Decisions include accepted, rejected, and abstained; verification
states distinguish verified agreement, evidence-grounded support, contradiction,
conflict, and unsupported status.

Gold applicability is likewise explicit. Each field--property slot is one of:
(i) applicable with a value, (ii) applicable but legitimately unknown, (iii) not
applicable, or (iv) outside an open-world gold scope. This prevents an abstention
from being silently scored as an incorrect accepted value and prevents missing gold
coverage from being misread as \NA.

\section{First-Principles Architecture Audit}

The audit separated essential from accidental complexity. Parser authority,
structured abstention, property-level provenance, and visible conflict states are
essential because they encode trust boundaries. Per-field prompt fragmentation,
manual grouping, repeated free-JSON/schema attempts, and redundant transformations
are historical until their value is demonstrated. Rigid output contracts remain
useful where they make failures observable; they are not justified merely as a way
to force a preferred answer.

The audit produced one conservative redesign rule: delete semantic orchestration
only after a replay-controlled experiment shows that a simpler approach is equal or
better. This avoids replacing old Model Capability Debt with new model dependence.

\section{The Four Fixed Architecture Variants}

\begin{table}[H]
\centering
\caption{Frozen semantic architecture variants. No additional variant is permitted in the study.}
\label{tab:variants}
\begin{tabularx}{\textwidth}{@{}p{0.8cm}p{3.0cm}Yp{2.3cm}@{}}
\toprule
ID & Name & Semantics & Semantic generation calls \\
\midrule
A & Deterministic only & Emits parser-owned schema and deterministic claims; no semantic model. & 0 \\
B & Dataset reasoner & Deterministic extraction plus at most one dataset-level semantic response for all unresolved targets. & $\leq 1$ per dataset \\
C & Replayed reasoner + verifier & Consumes the exact B response; checks evidence identity, field relevance, deterministic contradiction, and cross-claim conflict. & 0 new calls \\
D & Legacy semantic pipeline & Preserves per-field/manual-group scheduling, temperature 0.2, free-JSON then schema fallback, and the historical non-empty-evidence merge rule. & Variable, historically many \\
\bottomrule
\end{tabularx}
\end{table}

The B--C contrast is the cleanest intervention. A memoizing backend keys the full
request and provides C with the exact raw response and SHA-256 identity produced for
B. C is invalid if it generates a response, lacks a replay, or consumes a different
hash. Verification is deterministic and does not claim semantic entailment: an
existing, field-relevant citation establishes traceable support, not factual truth.

\section{Evaluation Methodology}

Let $V$ be applicable-value slots, $P$ accepted predictions on those slots, and
$K\subseteq P$ correct accepted predictions. We report
\begin{align}
  \text{end-to-end accuracy} &= \frac{\lvert K\rvert}{\lvert V\rvert},\\
  \text{coverage} &= \frac{\lvert P\rvert}{\lvert V\rvert},\\
  \text{selective risk} &= 1-\frac{\lvert K\rvert}{\lvert P\rvert}.
\end{align}
When $P$ is empty, selective accuracy/risk is undefined rather than forced to zero.
Abstentions lower end-to-end accuracy and coverage but do not enter the conditional
selective-risk denominator. Confidence is used only to order whole tie groups for a
right-continuous risk--coverage curve. AURC is integrated only over achieved
coverage; per-property curves are preferred when confidence scales differ.

Unsupported-claim rates are computed over model-proposed or model-accepted claims,
never diluted by deterministic claims. Cost accounting separates generated calls,
reused upstream calls, and effective architecture cost. A pairwise comparison is
non-comparable if either run is partial, replay identity fails, required gold is
missing, or a semantic target lies outside the scoring design.

\subsection{Adversarial evaluator validation}

The focused suites now contain \FocusedTests{} tests across the variant and backend
qualification modules, within \FullTests{} passing project tests. They cover
nonexistent, irrelevant, wrong-span, contradictory, and empty evidence; confidence
0 and 1; malformed and partial output; duplicates and conflicting claims; applicable
unknown and \NA; missing or mismatched replay; hidden retries; timeout; zero/all
accepted predictions; and confidence ties. Additional regressions protect failure
replay identity, per-call context accounting, legacy fallback telemetry, and partial
D provenance.

\begin{table}[H]
\centering
\caption{Core causal invariants enforced by code and tests.}
\begin{tabularx}{\textwidth}{@{}p{4.2cm}Y@{}}
\toprule
Invariant & Enforced behavior \\
\midrule
A call boundary & Always zero model calls. \\
B call boundary & At most one dataset-level generation call; hidden retries make the run partial. \\
C replay boundary & Zero generation calls and exact B raw-response hash, including malformed-response replay. \\
Evidence boundary & Missing, empty, cross-field, or wrong-span evidence cannot verify a claim. \\
Confidence boundary & Confidence never implies verification. \\
Applicability boundary & \NA, applicable unknown, applicable value, and outside-scope are distinct. \\
Comparability boundary & Partial execution, missing gold coverage, or replay failure suppresses causal deltas. \\
Legacy fidelity & Historical scheduling and merge behavior remain unchanged; all attempts are measured. \\
\bottomrule
\end{tabularx}
\end{table}

\section{Synthetic Development Experiment}

The development manifest contains \DevCases{} authored regression cases spanning
CSV, HDF5, and time-series structures. It is not a random sample and provides no
external-validity evidence. Only \DevSemanticCases{}/\DevCases{} cases contain any
dataset-level semantic target. One of those cases targets
\texttt{/campaign/meta/platform}, which is absent from gold; consequently every
formal A--B, B--C, and C--D comparison is non-comparable.

The local development backend was Qwen3.5 9B. The two actual B responses were
replayed byte-for-byte into C; the other seven cases correctly made no B/C call.
Table~\ref{tab:devresults} reports pooled descriptive diagnostics, not architecture
effects.

\begin{table}[H]
\centering
\caption{Nine-case synthetic development diagnostics. All B/C/D causal comparisons are non-comparable.}
\label{tab:devresults}
\scriptsize
\begin{tabularx}{\textwidth}{@{}lCCCCCC@{}}
\toprule
Variant & Comparable & End-to-end acc. & Coverage & Sel. risk & AURC & Calls \\
\midrule
A & yes & \AccA & \CovA & \RiskA & \AurcA & 0 \\
B & no  & \AccB & \CovB & \RiskB & \AurcB & \DevBCalls \\
C & no  & \AccC & \CovC & \RiskC & \AurcC & 0 new / \DevBCalls{} effective \\
D & no  & \AccD & \CovD & \RiskD & \AurcD & \DevDCalls \\
\bottomrule
\end{tabularx}
\end{table}

B proposed one model claim, accepted none, and added no scorable coverage. Its
effective cost was \DevBInput{} input tokens, \DevBOutput{} output tokens, and
\DevBLatency{} ms of model latency. C had the same effective upstream cost and zero
new generation. D proposed and accepted four claims, used \DevDCalls{} calls,
\DevDInput{}/\DevDOutput{} input/output tokens, and \DevDLatency{} ms. D recovered
two gold-covered logical values in one hard CSV case, but another change was
unscored. These facts do not establish superiority because the design is
non-comparable.

The Qwen3.5 backend abstained on a postal-code target even though its notes
recognized the semantics, and violated target scope on the unscored HDF5 case. A
stronger model might resolve either failure without architecture changes; this is a
hypothesis, not a result. The development run therefore yields no estimate of the
increment from A to B and no verification trade-off from B to C. The all-zero
observed B--C decision delta is uninformative, not evidence of no verifier effect.

\section{Backend Qualification}

Qualification asks whether a backend can execute the frozen contracts, not whether
its claims are correct. The fixed non-blind calibration manifest contains
\QualificationCases{} external datasets and \QualificationTargets{} semantic
targets. It is excluded from the future blind benchmark. B is run twice per dataset
at temperature 0 and seed 0; C immediately replays each response; D runs once with
its preserved policy.

\subsection{Invalidated 8K/thinking-on execution}

The first Qwen3.6-27B Q4\_K\_M execution used an 8,192-token loaded context while
thinking remained effectively enabled. Only 2/16 B attempts and 1/8 D dataset runs
were contract-valid; 12 B attempts terminated at the generation boundary, and a
10,044-token HDF5 prompt was rejected before generation. The run also exposed three
measurement defects: malformed B failures did not carry raw identity into C,
unparseable output could be counted target-scope valid, and partial D failures
dropped earlier-group/fallback provenance. The immutable artifact was invalidated;
none of its rates is used for backend selection.

The failure was informative as infrastructure diagnosis. LM Studio 0.4.18 documents
a fix for reasoning on/off being ignored by some custom configurations
\cite{lmstudio_0418}. We repaired only provenance and measurement, not prompts,
architectures, targets, or call budgets.

\subsection{Corrected 16K/thinking-off execution}

A one-call gate first established a contract-valid response with zero reasoning
tokens. The same Qwen3.6-27B Q4\_K\_M model \cite{lmstudio_qwen36} was then run with
thinking disabled and a \LoadedContext-token context. It passed every frozen
criterion (Table~\ref{tab:qualification}).

\begin{table}[H]
\centering
\caption{Corrected Qwen3.6-27B operational qualification.}
\label{tab:qualification}
\begin{tabularx}{\textwidth}{@{}Yrrc@{}}
\toprule
Criterion & Observed & Required & Result \\
\midrule
Semantic-opportunity datasets & 8 & $\geq 8$ & pass \\
B contract-valid rate & 1.0 & $\geq 0.95$ & pass \\
B target-scope-valid rate & 1.0 & 1.0 & pass \\
B/C replay identity & 1.0 & 1.0 & pass \\
B exact repeatability & 1.0 & 1.0 & pass \\
B timeout / truncation rates & 0.0 / 0.0 & 0.0 / 0.0 & pass \\
B context fit / observed telemetry & 1.0 / 1.0 & 1.0 / 1.0 & pass \\
D contract-valid rate & 1.0 & $\geq 0.95$ & pass \\
D timeout / truncation rates & 0.0 / 0.0 & 0.0 / 0.0 & pass \\
D context fit / observed telemetry & 1.0 / 1.0 & 1.0 / 1.0 & pass \\
\bottomrule
\end{tabularx}
\end{table}

All 16 B attempts ended with JSON and \texttt{finish\_reason=stop}; repeats were
byte-identical for all eight datasets. Every C input hash matched B and every C
semantic generation count was zero. The maximum observed B context was
\QualMaxObservedContext{} tokens; the maximum requested
\texttt{input+max\_tokens} budget was \QualMaxBudgetContext{}, below the loaded
limit.

\begin{table}[H]
\centering
\caption{Descriptive operational cost in the corrected qualification.}
\small
\begin{tabularx}{\textwidth}{@{}lRRRY@{}}
\toprule
Condition & Calls & Input tokens & Output tokens & Mean latency \\
\midrule
B (two repeats) & \QualBCalls & \QualBInput & \QualBOutput & \QualBMeanLatency{} s per attempt \\
C generation & 0 & 0 new & 0 new & deterministic replay/verification \\
D (one run) & \QualDCalls & \QualDInput & \QualDOutput & \QualDMeanLatency{} s per dataset \\
\bottomrule
\end{tabularx}
\end{table}

D's 59 calls are an observed consequence of historical scheduling: the HDF5 case
alone required 22 field calls. One \texttt{functional\_traits.csv} free-JSON shape
attempt failed and the preserved schema fallback succeeded. The cost gap is real;
an accuracy--cost trade-off has not yet been estimated.

\section{What the Current Results Do and Do Not Establish}

\paragraph{Established facts.}
The harness enforces the four fixed architectures, exact B/C replay, zero C
generation, fail-closed contracts, explicit applicability, and per-call telemetry.
The Qwen3.6-27B Q4\_K\_M condition is operationally eligible under the corrected
runtime. The legacy architecture has a substantially larger call surface on the
calibration tasks.

\paragraph{Not established.}
The study has no blind effect estimate. The nine-case development set is too sparse
and incompletely gold-covered to estimate A--B, B--C, or C--D effects. Backend
qualification contains no gold. It cannot show that dataset-level reasoning helps,
that verification lowers risk, that D is more accurate, or that Qwen generalizes.

\paragraph{Model selection.}
There is no operational reason to replace Qwen3.6-27B with GLM, Llama, or a smaller
Qwen. Alternative families may be predeclared sensitivity conditions, but adding
them post hoc to rescue an unfavorable architecture result would confound backend
and architecture selection. Model cards and exact runtime identity should be
reported as first-class artifacts \cite{mitchell2019modelcards}.

\section{Blind External Evaluation Protocol}

Two annotators who did not develop the system will independently label the blind
gold. They will distinguish applicable value, applicable unknown, \NA, and
outside-scope slots, then adjudicate disagreements without observing architecture
outputs. Dataset-level paired effects, not pooled slots alone, will be the primary
analysis unit. The development set will not be used as evidence of generalization.
Dataset documentation follows the broader principle that collection, intended use,
and limitations should be explicit \cite{gebru2021datasheets}.

The primary contrasts are predeclared:
\begin{itemize}
  \item A--B: added correct/incorrect claims, coverage, selective risk, unsupported claims, calls, tokens, and latency;
  \item B--C: B-accepted claims rejected or abstained by C, incorrect claims removed, correct claims lost, unsupported claims removed, and the risk--coverage trade-off;
  \item C--D: quality, selective metrics, unsupported claims, calls, tokens, latency, and failure surface, only when D is comparable.
\end{itemize}

\subsection{Blind-freeze provenance placeholders}

\begin{center}
\fcolorbox{FreezeRed}{FreezeBg}{%
\begin{minipage}{0.93\linewidth}
\textbf{STOP: the fields below are intentionally unresolved.}\par
They must be content-addressed and countersigned before blind gold is unsealed.
Leaving a field blank blocks interpretation; it must never be silently replaced by
the development or qualification artifact.\medskip

\begin{tabularx}{\linewidth}{@{}>{\raggedright\arraybackslash}p{5.2cm}Y@{}}
Blind sampling registration receipt & \placeholder \\
Blind candidate-frame SHA-256 & \placeholder \\
Blind manifest SHA-256 & \placeholder \\
Independent annotation package SHA-256 & \placeholder \\
Adjudicated gold package SHA-256 & \placeholder \\
Annotation vocabulary/handbook SHA-256 & \placeholder \\
Annotator pseudonyms and independence attestations & \placeholder \\
Calibration-gate result and agreement statistics & \placeholder \\
Frozen backend-registry SHA-256 & \placeholder \\
Exact Qwen GGUF checkpoint SHA-256 & \placeholder \\
LM Studio 0.4.18 runtime/version receipt SHA-256 & \placeholder \\
Frozen source commit/tree SHA-256 & \placeholder \\
Frozen analysis-plan SHA-256 & \placeholder \\
Prompt/schema bundle SHA-256 & \placeholder \\
Gold unsealing timestamp and countersignatures & \placeholder \\
\end{tabularx}
\end{minipage}}
\end{center}

The known pre-freeze contract hashes are reported in
Appendix~\ref{app:artifacts}; they do not substitute for the final registry,
checkpoint, source-tree, or gold receipts.

\section{Threats to Validity}

\textbf{Construct validity.} Evidence identity and field relevance do not prove
semantic entailment. C should be interpreted as deterministic trust-boundary
verification, not an oracle. Confidence values may be incomparable across
properties; per-property selective curves are therefore necessary.

\textbf{Internal validity.} B/C replay is tightly controlled, but D intentionally
retains temperature 0.2 and historical fallback behavior. Backend state, runtime
version, context, checkpoint, and prompt hashes must remain frozen. Any partial
response or missing replay suppresses the corresponding comparison.

\textbf{External validity.} The nine development cases are synthetic and the eight
qualification cases are a fixed operational calibration set already exposed to
development. Neither is a population sample. Only the future independently
annotated corpus can support generalization claims.

\textbf{Statistical conclusion validity.} The blind analysis must use paired
dataset-level effects and confidence intervals. A clean null result is acceptable.
The current manuscript reports descriptive pooled metrics only to document harness
behavior.

\textbf{Reproducibility.} The exact checkpoint-file hash and machine-observed LM
Studio version are still missing. The version is operator-reported. These are
explicit freeze blockers rather than hidden caveats.

\section{Conclusion}

The durable system objective remains dataset input to deterministic schema output,
with optional evidence-constrained semantic enrichment. The present work reduces
the semantic architecture question to four auditable variants and validates the
measurement boundaries needed for causal interpretation. Development data do not
rank the variants. Corrected backend qualification shows that one dataset-level
Qwen3.6-27B response can be generated repeatably and replayed into deterministic
verification, while the legacy condition exposes a much larger call surface.

The central research answer is intentionally deferred: whether semantic reasoning
adds useful claims, and whether deterministic verification removes more errors than
correct claims, can only be learned after the blind provenance block is completed
and independent gold is unsealed. Until then, the scientifically correct result is
a validated protocol, an operationally eligible backend, and no architecture winner.

\appendix
\section{Content-Addressed Current Artifacts}
\label{app:artifacts}

\begin{longtable}{@{}p{5.2cm}p{10.0cm}@{}}
\toprule
Artifact & SHA-256 \\
\midrule
\endhead
Development manifest & \artifacthash{\DevManifestHash} \\
Development report & \artifacthash{\DevReportHash} \\
Qualification manifest & \artifacthash{\QualManifestHash} \\
Controlled vocabulary & \artifacthash{\VocabularyHash} \\
Invalidated 8K qualification & \artifacthash{\InvalidQualHash} \\
Corrected qualification & \artifacthash{\ValidQualHash} \\
Qualification assessment & \artifacthash{\QualAssessmentHash} \\
Dataset reasoner prompt & \artifacthash{\DatasetPromptHash} \\
Dataset response schema & \artifacthash{\DatasetSchemaHash} \\
Legacy prompt & \artifacthash{\LegacyPromptHash} \\
Legacy response schema & \artifacthash{\LegacySchemaHash} \\
\bottomrule
\end{longtable}

\section{Current Reproduction Status}

At manuscript generation, \FullTests{} project tests passed. The relevant Ruff
checks passed, all 316 JSON artifacts parsed successfully, and
\texttt{git diff --check} passed. Seven pre-existing full-tree lint findings in
unrelated legacy files remain documented and were not modified. The two PDFs were
rendered from these TeX sources using Tectonic 0.16.9 \cite{tectonic0169}.

\bibliographystyle{plain}
\bibliography{references}

\end{document}
